\section{Minimal rank-one constant completion}
\label{sec:constant-completion}

An unsafe orthogonal projection has a canonical repair.  It adds
exactly the direction in which the constant vector is missing.

\begin{theorem}[Minimal constant completion]
\label{thm:minimal-completion}
Let \(P\) be an orthogonal projection and let
\[
 q=(I-P)\one.
\]
If \(q=0\), put \(P^\sharp=P\).  If \(q\ne0\), define
\begin{equation}
 P^\sharp x
 =
 Px+
 \frac{\langle x,q\rangle}{\|q\|_2^2}q.
\label{eq:completed-projection}
\end{equation}
Then:
\begin{enumerate}[label=(\roman*),leftmargin=2.2em]
\item \(P^\sharp\) is the orthogonal projection onto
  \[
  \Ran P\oplus\C q
  =
  \Ran P+\C\one;
  \]
\item \(P^\sharp\one=\one\);
\item \(\Ran P^\sharp\) is the smallest subspace containing
  \(\Ran P\) and \(\one\);
\item for every \(A\),
  \begin{equation}
  D(P^\sharp A)
  =
  D(PA)+
  \frac{|\langle A,q\rangle|^2}{\|q\|_2^2},
  \label{eq:completion-energy}
  \end{equation}
  and
  \begin{equation}
  Z(P^\sharp A)=Z(A).
  \label{eq:completion-zero}
  \end{equation}
\end{enumerate}
\end{theorem}

\begin{proof}
Because \(q\in\Ker P=(\Ran P)^\perp\), the second term in
\eqref{eq:completed-projection} is the orthogonal projection onto
\(\C q\), perpendicular to \(P\).  Their sum is therefore the
orthogonal projection onto \(\Ran P\oplus\C q\).

Since \(\one=P\one+q\),
\[
 \Ran P+\C\one
 =
 \Ran P\oplus\C q.
\]
Also
\[
 \langle\one,q\rangle
 =
 \langle P\one+q,q\rangle
 =
 \|q\|_2^2,
\]
so \eqref{eq:completed-projection} gives
\[
 P^\sharp\one=P\one+q=\one.
\]
Any subspace containing both \(\Ran P\) and \(\one\) contains
\(q=\one-P\one\), proving minimality.  Orthogonality of the two
summands in \(P^\sharp A\) gives
\eqref{eq:completion-energy}, and
\eqref{eq:completion-zero} follows from
\cref{prop:preservation-criterion}.
\end{proof}

\begin{corollary}[Repaired monotonicity]
\label{cor:repaired-monotonicity}
Let \(A\ne0\).  If \(P^\sharp A=0\), then \(Z(A)=0\).  Otherwise,
\[
 \rho(A)\le\rho(P^\sharp A),
\qquad
 \Flat(A)\le\Flat(P^\sharp A).
\]
\end{corollary}

\begin{proof}
Apply \cref{thm:safe-projection-law} to \(P^\sharp\).
\end{proof}

\begin{remark}[The repair is exact but not free]
\label{rem:repair-not-free}
The added coefficient is
\[
 \langle A,q\rangle
 =
 Z(A)-Z(PA).
\]
Thus rank-one completion identifies the unique missing scalar, but
does not estimate it.  If that scalar is already available from an
exact provenance ledger, the completion restores a safe
same-frame projection.  If it is the unknown arithmetic zero mode,
writing \(P^\sharp\) does not solve the arithmetic problem.
\end{remark}

\begin{proposition}[Uniqueness at minimal range]
\label{prop:minimal-uniqueness}
The operator \(P^\sharp\) is the unique orthogonal projection \(Q\)
such that
\[
 \Ran Q=\Ran P+\C\one.
\]
If \(q\ne0\), every constant-preserving orthogonal projection whose
range contains \(\Ran P\) has rank at least
\(\operatorname{rank}P+1\), with equality precisely when its range
equals \(\Ran P^\sharp\).
\end{proposition}

\begin{proof}
An orthogonal projection is uniquely determined by its range.
When \(q\ne0\), \(\one\notin\Ran P\), so adjoining \(\one\) raises
the dimension by one.  Every range containing both objects contains
their span.
\end{proof}
